% !TEX root = ../main.tex
\chapter{Governed Retrieval: A Reference Architecture}
\label{ch:governed-rag}
\label{sec:governed-rag}

Section~\ref{sec:measurability} established a requirement rather than a preference:
the grounding rate is measurable only where the evidential base is explicit at
inference time. Section~\ref{sec:retrieval-insufficient} established that standard
retrieval supplies one of seven grounding dimensions reliably and none of the three
that are relations to $c$. This section sets out what closes the gap.

The distinction I will use throughout:

\begin{quote}
Standard retrieval retrieves evidence. \emph{Governed} retrieval establishes the
controlled conditions under which retrieved evidence may support an institutional
assertion.
\end{quote}

This is architecture, not prompt engineering. None of the eight components below can
be implemented as an instruction to the model, and Section~\ref{sec:missing} says
why: instructions enter as tokens and are processed by the same mechanism as the
content, so an instruction to abstain is a probabilistic influence on a
continuation rather than a precondition on output.

\section{The underspecified-request principle}
\label{sec:holodeck}

Before the components, the requirement that motivates the first two of them.

\begin{quote}
Natural language is an incomplete specification for consequential action.
\end{quote}

A fluent request routinely leaves unresolved: the referents; the jurisdiction; the
effective time; the applicable policy; the objective; the requester's authority; the
constraints; the side effects; the reversibility; the threshold at which action
should be taken; and whether the requester wants analysis, a recommendation,
execution, or a binding determination.

The point is not that natural language is a poor interface. It is an excellent
interface, which is why these systems are worth building. The point is that a
fluent request cannot be treated as a \emph{complete operational specification} in a
regulated decision or execution workflow, and that the gap between the two is
invisible in the request itself. A well-formed sentence does not announce which of
its parameters were never bound.

Two consequences. The system must force or validate material bindings before
crossing into consequential action---which is I10. And the boundary at which it
crosses has to be an explicit architectural decision, because by
Section~\ref{sec:rlhf} the training regime supplies pressure to answer rather than
to ask.

\section{Eight components}

\paragraph{1. Authoritative corpus control.} Sources are approved, authenticated,
and governed. Each carries stable identity, version, effective date, jurisdiction,
authority level, and applicability metadata, and sources can be pinned to a corpus
release or a decision-time snapshot. This supplies I8, I9, and the substrate for
I11 and I14. It is the component most familiar from ordinary data governance and,
in my experience, the one most often assumed to be sufficient on its own.

\paragraph{2. Context resolution before assertion.} Material variables---$j$, $t$,
$f$, $u$---are resolved, collected, or explicitly marked unknown before an
assertion is produced. Where material context is missing the system requests it,
qualifies the output, declines to produce a decision-ready assertion, or escalates.
This supplies I10 and is the only defense against illicit universalization
(Section~\ref{sec:illicit}). It is also the component with the highest product cost,
because it makes the interaction less immediately satisfying, and it is therefore
the first to be cut.

\paragraph{3. Metadata-aware retrieval.} Retrieval filters or ranks on scope
constraints, not on semantic similarity alone: jurisdiction predicates, effective
date windows, population or product class, decision type. The retrieval event and
the resulting evidence set are retained. This supplies I11, and it is the direct
answer to the observation of Section~\ref{sec:retrieval-insufficient} that a
paediatric and an adult dosing passage are near-neighbours in embedding space.

\paragraph{4. Claim decomposition.} Complex outputs are decomposed, where feasible,
into atomic claims, each with linked evidence and an explicit support relation. This
supplies I12. A response-level citation list does not permit checking the third
sentence, and the third sentence is frequently the one relied upon.

\paragraph{5. Evidence-to-assertion verification.} Citation presence is
insufficient. The system checks whether the cited evidence supports the claim, and
whether it does so under $c$ and at the strength asserted
(Section~\ref{sec:entails-supports}). Controls may combine deterministic rules,
structured policy logic, source checks, independent model review, and human
validation, proportioned to risk. This supplies I13 and is where the
$\Sensitive$ condition of Definition~\ref{def:grounding} is actually enforced
rather than assumed.

\paragraph{6. Output gating.} If evidence is missing, stale, conflicting,
inapplicable, or insufficient for the strength of claim requested, the system does
not present an unqualified assertion as decision-ready. It returns bounded
information, an explicit uncertainty, a request for missing context, or an
escalation path. This supplies I17, and it must sit outside $M$.

\paragraph{7. Audit record and replay.} Retain the prompt or structured input; the
resolved context; source identifiers and versions; retrieved passages; model
version; prompt or template version; system configuration; verification results;
human interventions and approvals; the final response; and timestamps with
retention identifiers. This supplies I14 and I16, and it is what answers the
question posed at the start of Section~\ref{sec:intro}.

\paragraph{8. Authorization and tiered reliance.} Permitted use depends on the
consequence of reliance, not only on model quality. Higher-consequence uses require
stronger constraints, independent validation, and where appropriate a qualified
human decision-maker with the standing to be accountable. This supplies I15, and the
tiers of reliance that make it proportionate are developed with the definition of
warrant.

\section{The flow}

\begin{center}
raw generative output
$\rightarrow$ context resolution
$\rightarrow$ controlled retrieval
$\rightarrow$ applicability filtering \\
$\rightarrow$ evidence-to-claim verification
$\rightarrow$ authorization or escalation
$\rightarrow$ auditable decision support
\end{center}

Two things about this sequence carry most of its value, and neither is the ordering.

Each arrow is a \emph{stopping point}. Context resolution can return a request for
missing bindings; applicability filtering can exclude every candidate source; verification
can find that the evidence does not reach the claim; the defeater search can find contrary
authority; and authorization can decline. A pipeline in which the only outcome is an
answer has implemented the boxes without the arrows, and it is the arrows that make the
conditions enforceable rather than advisory.

Each arrow is also a \emph{record}. What the audit trail contains is not a summary of the
final answer but the sequence of determinations that produced it: which context was
resolved and how, which sources survived filtering and which were excluded and why, what
the verification returned, what the defeater search returned, and who authorized the
result. That is what makes the flow reconstructable months later by someone who was not
present, which is the requirement Section~\ref{sec:intro} opened with.

The sequence is not a confidence gradient. It is a sequence of distinct conditions, and an
assertion that satisfies six of them is not ninety per cent warranted; it is unwarranted
for a reason that the record should name.

\section{What each intervention buys}

Common interventions address different invariants, and conflating them is how programs
spend a great deal and move one number. The table below is best read as a purchasing
guide: the left column is what a program funds, the middle column is what that
purchase actually delivers against Table~\ref{tab:invariants-system}, and the right
column is what it is frequently assumed to deliver and does not.

Two features of it are worth pointing out before reading it. Retrieval appears with two
distinct entries in the middle column because it does two unrelated things: it reduces
fabrication, which is why it is usually funded, and it supplies an identified evidential
base, which is what makes \eqref{eq:rate} measurable at all and is the more durable of
the two. And no single row delivers more than a handful of invariants, which is the
practical form of an argument about ownership: there is no purchase
that makes a system auditable, only a set of purchases that between them cover the
conditions.

\begin{table}[ht]
  \centering
  \small
  \begin{tabular}{@{}>{\raggedright\arraybackslash}p{3.3cm}>{\raggedright\arraybackslash}p{4.6cm}>{\raggedright\arraybackslash}p{5.4cm}@{}}
    \toprule
    Intervention & Buys & Does not buy \\
    \midrule
    Standard retrieval
      & I8; makes \eqref{eq:rate} measurable; reduces \textsc{Fabrication}
      & I10--I11, I13; form invariance; applicability \\
    \addlinespace
    Better base model
      & fewer fabrications on covered topics
      & I8--I19; may worsen \textsc{Lucky} by guessing better \\
    \addlinespace
    Verifier loops
      & I13 on checkable claims; some of I2--I4
      & grounding, unless the verifier checks evidence rather than form \\
    \addlinespace
    Persistent memory
      & a precondition for I7 and I14
      & either of them---storage is not retrieval discipline \\
    \addlinespace
    Context resolution
      & I10, and the substrate for I11
      & anything on the generation path \\
    \addlinespace
    Output gating
      & I17, and thereby enforceability of the rest
      & correctness of the evidence itself \\
    \bottomrule
  \end{tabular}
  \caption{Retrieval does double duty in most deployments: it is funded to improve
    accuracy, and its more durable contribution is making the grounding rate
    observable at all.}
  \label{tab:interventions}
\end{table}

The row worth dwelling on is the second. A better base model improves the numbers
most programs report while leaving every system invariant untouched, and its
effect on the lucky cell runs the wrong way: a model that guesses better produces
more cell-(ii) assertions and scores higher for doing so.

\section{Structured evidence is the strongest case}
\label{sec:structured}

The components above are described in terms of retrieved passages, which
understates what is available. The highest-assurance form of $E$ in an enterprise is
not a document at all. It is a \emph{query against a system of record}, with the query
text, the parameters, the result set, and the as-of state all recorded.

Structured evidence satisfies the grounding conditions more cleanly than text
retrieval does, because the evidence has a schema. Authenticity follows from the source
system's own controls. Versioning is the table's temporal model rather than a document
revision. Applicability becomes a set of predicates over columns---jurisdiction,
effective date range, product class, population---evaluated rather than judged.
Sufficiency becomes a completeness check over required fields. Traceability is the
recorded query. And the defeater check of \eqref{eq:nodefeater} becomes a query for
contradicting rows rather than a search for contradicting prose.

The practical consequence is a preference order that most retrieval-first designs
invert: where a claim can be answered from a governed table, it should be, and the
model's role is to render the result rather than to find it. Text retrieval is the
fallback for claims whose evidence exists only as prose, and it inherits every
judgment that the schema would otherwise have decided.

\section{What this costs}
\label{sec:cost}

The controls are not free and a proposal that does not say so is not usable. The
multipliers can be given structurally, without benchmarks, because they follow from how
many model calls and retrievals each control requires.

Let a response contain $m$ consequential claims, each with up to $k$ candidate
supporting sources. Generation is one call. Claim decomposition is one call plus $O(m)$.
Support verification is $O(mk)$ calls where it is model-based, and $O(1)$ per claim where
the check is deterministic. The defeater search is one additional retrieval per claim,
$O(m)$. Coverage checking is $O(m)$ against an enumerated rule and unbounded against a
rule that must be read out of prose. So the dominant term is $mk$, and the design levers
follow directly: reduce claims per response, reduce candidate sources per claim, or
replace model verification with deterministic checks. A ten-claim response over three
candidate sources each is on the order of thirty verification calls against one
generation call---which is the difference between a system that costs what a chat costs
and one that does not.

Storage grows with retained evidence rather than with traffic, and
Section~\ref{sec:retention} constrains what may lawfully be kept.

\begin{table}[ht]
  \centering
  \small
  \begin{tabular}{@{}>{\raggedright\arraybackslash}p{3.9cm}>{\raggedright\arraybackslash}p{3.1cm}>{\raggedright\arraybackslash}p{6.6cm}@{}}
    \toprule
    Component & Cost profile & Notes \\
    \midrule
    Corpus control, metadata
      & one-off plus upkeep
      & Ontology and curation effort; no per-request cost \\
    \addlinespace
    Metadata-aware retrieval
      & negligible per request
      & Predicate filtering is cheaper than the embedding search it constrains \\
    \addlinespace
    Context resolution
      & one round trip, or a form
      & Latency and interaction cost, not compute \\
    \addlinespace
    Claim decomposition
      & multiplies calls
      & Linear in claims per response \\
    \addlinespace
    Support verification
      & multiplies again
      & The dominant cost. Deterministic checks are cheap; model-based
        entailment is per claim per source; human adjudication is the expensive tail \\
    \addlinespace
    Defeater search (I19)
      & one extra retrieval
      & The governed contrary query, run under the same applicability predicates.
        Cheap, and it needs the precedence order to exist \\
    \addlinespace
    Coverage checking (I18)
      & per claim
      & Cheap where the rule's conditions are enumerated; expensive where they must be
        inferred from prose \\
    \addlinespace
    Audit record
      & storage, growing
      & See Section~\ref{sec:retention} on what may lawfully be kept \\
    \bottomrule
  \end{tabular}
  \caption{Cost profile of the controls. The expensive items are per-claim
    verification and human adjudication, which is precisely why tiering
    makes them proportional to consequence.}
  \label{tab:cost}
\end{table}

Tiered reliance exists for this reason and not only for
governance tidiness. If every request bore the full control set the economics would
not work, and a design that applies Tier~3 controls to Tier~1 traffic will be
abandoned rather than descoped.

\section{Retention against minimization}
\label{sec:retention}

The audit record of component~7 conflicts with another control regime, and the
conflict should be stated rather than discovered in review. Prompts, resolved contexts,
and retrieved passages routinely contain personal or special-category data. Data
minimization and erasure obligations cut directly against the retention that
reconstructability requires.

Four resolutions, none complete on its own. Retain references, identifiers, and hashes
in place of content wherever the content is recoverable from a governed source, so the
audit record points at evidence rather than copying it. Segregate the audit store and
control access to it separately from the application. Align retention to the retention
schedule of the \emph{decision} the assertion supported, rather than choosing a new
period for the AI system. And where erasure is exercised, preserve the structure of the
record---that a claim was made, on what evidence class, under what authority---while
removing the identifying content.

The point for this book is narrower than the resolutions. Reconstructability is a
requirement in tension with minimization, the tension is a design decision rather than
an implementation detail, and it requires a fifth function: privacy and records
management, which joins the three that own the rest.

\section{What this does not claim}

Governed retrieval does not make these systems infallible, and nothing above should
be read as proposing that it does. Cell (iii) survives it entirely: a controlled,
versioned, applicable, authoritative corpus can still be wrong, and no amount of
architecture converts a bad source into a good one. Verification is relative to the
checks implemented and the assumptions they encode. Human review introduces its own
failure modes and its own capacity limits. And the floor of
Section~\ref{sec:floor} does not move.

The function of the architecture is narrower and worth stating precisely. It
transforms otherwise unlicensed generated text into a controlled, inspectable input
to an accountable institutional process---one whose failures are attributable,
whose evidence is identifiable, and whose reasoning can be reconstructed after the
fact by someone who was not present. That is a lower bar than reliability and a
much higher bar than accuracy.

\section{Bounding the Request Space}
\label{ch:bounded-requests}
\label{sec:bounded}

Section~\ref{sec:canonical-form} argued that the absence of a canonical form is what
makes accuracy uncertifiable, and Section~\ref{sec:normalizer} argued that training a
semantic normalizer is not the remedy. There is a third option, and it is the one an
architect reaches for first: if the model will not canonicalize its input and cannot
be taught to, \emph{do not give it free-form input at all}.

Restricting the interface is not a workaround. It is the direct application of the
diagnosis: canonicalization moves upstream of the model, into a schema, where it is
inspectable and versioned. This section works through what that buys, what it does
not, and where the residual risk goes---because it does not disappear.

\subsection{The construction}

Replace free text with a typed request. Instead of ``is this approved for children?''
the interface accepts an object:

\begin{center}
\small
\begin{tabular}{@{}l>{\raggedright\arraybackslash}p{7.2cm}@{}}
\toprule
\texttt{intent} & \texttt{INDICATION\_CHECK} (from a closed set) \\
\texttt{subject\_id} & a key into a system of record \\
\texttt{indication} & a coded value \\
\texttt{jurisdiction} & required \\
\texttt{as\_of} & required \\
\texttt{population} & required, from a controlled vocabulary \\
\bottomrule
\end{tabular}
\end{center}

The object is rendered into a fixed template by code the organization controls and
versions. Requests are stateless: one request, one fresh context.

The fields are not arbitrary, and the discipline that selects them is the point rather
than the example. Every component of $c$ appears, because
Section~\ref{sec:resolved-context} established that a claim whose context is unresolved is
not yet evaluable---so jurisdiction, valid time, and intended use are required fields, and
a request missing any of them does not validate. The material facts $f$ enter as keys into
systems of record rather than as free text, so that the entity under discussion is
identified rather than described. The intent is drawn from a closed set, because the
decision class determines which sources have standing under $\standing$ and therefore what
counts as authoritative. And coded values replace prose wherever a controlled vocabulary
exists, since that is what allows applicability to be evaluated as a predicate in
Section~\ref{sec:governed-rag}'s third component rather than judged.

What the schema encodes, in other words, is the set of bindings that would otherwise have
to be guessed. Anything the type system requires is a thing the system can no longer
silently assume.

\subsection{What this fixes}

Against Table~\ref{tab:invariants-system}, the gains are specific and substantial.

\paragraph{I1 becomes satisfied by construction.} Not because the model acquired form
invariance, but because the opportunity to violate it has been removed. Two
materially equivalent requests are now \emph{the same request}: the equivalence class
of Section~\ref{sec:canonical-form} has exactly one member, and the canonical form is
supplied by the schema.

\paragraph{The request space becomes enumerable.} It is a product of typed fields
rather than $\mathcal{V}^{*}$. It may be combinatorially large, but it is
\emph{stratifiable}: coverage can be designed by jurisdiction, date band, population,
and intent, with known sampling properties. That is ordinary test design. The
validation problem becomes well posed in the sense of
Section~\ref{sec:not-easier}---the same sense in which it is well posed for a
classifier.

\paragraph{I10 is enforced by the type system rather than detected.} A request without
a jurisdiction does not validate. The misapplied category of
Section~\ref{sec:cell-v} does not need to be caught downstream, because the binding
it depends on is a precondition of submission. This is the single largest gain, and
it is available with no machine learning at all.

\paragraph{I7 becomes achievable.} Statelessness removes context accumulation, which
is the dominant source of within-session drift. Combined with the reproducibility
conditions of Section~\ref{sec:determinism}, replay becomes a realistic control
rather than an aspiration.

\subsection{What this does not fix}

\paragraph{The output space is still $\mathcal{V}^{*}$.} Bounding the input bounds the
input. The model can still emit any string, so ungrounded assertion survives
untouched: a confident figure resting on nothing is exactly as available as before.
Grounding is I13 and no amount of input discipline supplies it.

\paragraph{Errors are still not label mismatches.} If the response is prose, it can be
partly right, right for the wrong reason, or adequate for a weaker claim than the one
made. The atomic countable error of the classifier case does not follow from a typed
input.

\subsection{Where the problem goes}
\label{sec:relocation}

This is the part that decides whether the design is an improvement or a disguise.
\textbf{Who fills in the request?}

If a person translates their need into the typed object, the ambiguity has not been
eliminated; it has been moved into that translation, and it is unrecorded unless the
system captures it.

If a language model performs the translation---which is what is usually built,
because it preserves the interface that made the technology attractive---then the
unbounded input space has been reintroduced at the parsing stage. Form-invariance
failure now determines \emph{which request object is constructed}. A paraphrase can
produce a different jurisdiction, a different population, or a different intent, and
everything downstream will process that object correctly and produce a well-grounded
answer to the wrong question.

This is worse than the original failure in one specific respect: it is less visible.
An ungrounded assertion at least looks like an assertion, and the audit record shows
what evidence was and was not used. A misparsed request produces a record that is
internally consistent and complete. Every control passes. The evidence applies to the
context in the record, and the context in the record is not the user's context.

The mitigation is not technical. It is to \textbf{render the resolved request back and
require confirmation before proceeding}, which converts a silent failure into a
visible one and puts the binding decision where accountability already sits. This is
the pattern enterprise systems have always used before consequential execution, and
the argument of this book is that a natural-language front end does not earn an
exemption from it. Where confirmation is impractical, the parse should be treated as
part of the assertion for audit purposes: recorded, versioned, and measured for
accuracy against reviewed samples like any other component.

\subsection{Constraining the output, and what that does and does not buy}
\label{sec:constrained-output}

The natural complement is to constrain the output as well: grammar-constrained
decoding into a schema, with an enumerated verdict field, required evidence
references, and an explicit \texttt{INSUFFICIENT\_EVIDENCE} variant.

This is worth doing and its benefits should be stated exactly, because it is easy to
credit it with more than it delivers.

It \emph{does} guarantee that the output parses, that the output space is enumerable,
and that errors become atomic and countable again---restoring on the output side the
well-posedness that typed requests restore on the input side. It \emph{does} make
abstention a first-class value rather than a hoped-for phrasing, and it makes
per-field evidence linkage (I12) structurally enforceable rather than conventional.

It does \emph{not} satisfy I17. A grammar constrains which strings can be emitted; it
does not determine which the model selects. The choice between a substantive verdict
and the abstention variant is made by the same next-token computation as everything
else, subject to the same pressures identified in Section~\ref{sec:rlhf}. Constrained decoding makes abstention \emph{expressible}, not
\emph{correctly triggered}.

I17 is therefore satisfied by the check, not by the schema. The evidence conditions of
Definition~\ref{def:grounded-evidence} are evaluated outside the model, and the gate
must be able to \emph{override} a substantive verdict the model produced---replacing it
with a qualification, a refusal, or an escalation---when those conditions are not met.
A system that trusts the model's own choice of variant has implemented a vocabulary
for abstention and no gate.

\subsection{Where this is practical}

The design suits high-volume, repetitive, well-typed decisions: eligibility
determinations, coding against a rulebook, standard disclosures, reconciliations,
completeness checks. The user is an operator working a queue against a known decision
class, and the intent set is small and stable.

It suits exploratory work badly, and that matters, because exploration is a large part
of why these systems were adopted. An analyst pursuing an unanticipated question has
no intent to select.

Two costs deserve naming.

\paragraph{The intent taxonomy is an ontology project.} Enumerating decision classes,
coding their fields, and maintaining controlled vocabularies is the same work as
building a conformed dimensional model, with the same effort profile and the same
tendency to be underestimated.

\paragraph{An incomplete taxonomy fails quietly.} A question outside the intent set
does not raise an error. It is routed to the nearest available intent, which is a
misparse of the kind described in Section~\ref{sec:relocation}. Twenty intents may
cover most of the volume, and the residual tail is disproportionately where the
unusual, high-consequence cases live. A closed intent set therefore needs an explicit
\emph{no matching intent} outcome and a route for it, or the constraint will silently
manufacture the failure it was introduced to prevent.

\subsection{The effect on what the system is}

One consequence should be faced directly, because it changes how the component should
be described and governed.

If the input is typed, the output is schema-constrained, and the evidence conditions
are checked outside the model, then it is fair to ask what the model is contributing.
Frequently the honest answer is: retrieval, summarization, and the rendering of a
determination into readable language---while the determination itself is made by
deterministic rules over structured evidence.

That is a legitimate and valuable architecture, and it is one enterprises already know
how to validate. But it \emph{reclassifies} the system. The model is a retrieval and
presentation layer over a decision service, not the decision-maker, and the assurance
argument should say so. A program that describes such a system as ``AI-driven
decisioning'' is overstating the model's role and, more importantly, misdirecting its
own validation effort toward the component that is doing the least consequential work.

There is a further consequence worth facing, because it follows from combining this
section with Section~\ref{sec:structured} and a reader will combine them. If structured
evidence is the strongest available form of $E$, and typed requests plus constrained
output are what reach the higher reliance tiers, then the architecture that best
satisfies this book's conditions is one in which the model does comparatively little of
the deciding. Taken to its conclusion, that is an argument for routing decidable claims
\emph{away} from the model altogether.

I take that to be a feature rather than an embarrassment. A framework that identifies
which decisions do not need a language model is more useful than one that assumes every
decision does. What remains for the model, once the decidable residue is routed away, is
the genuinely open-textured part: synthesis across heterogeneous sources, questions whose
evidence exists only as prose, interpretation where the governing text is ambiguous, and
interaction with people who cannot state their need in a schema. That is a smaller claim
about the technology than the field usually makes, and a more defensible one---and it is
the part where the measurement problem of Section~\ref{sec:auditable} genuinely bites,
because it is the part no schema can decide.

Tiered reliance is where this resolves. Bounding the request
space is not a general answer to the problem of ungrounded assertion; it is the
mechanism by which a narrow, well-understood decision class can be moved to a higher
reliance tier than free-form interaction can support. Free-form interfaces are
appropriate at Tiers~0 to~2. Reaching Tier~3 and above is, in practice, an exercise in
bounding what may be asked and what may be answered.
